\subsection{DP de Intervalos (Corte Óptimo de una Barra)}
\label{alg:6-8}

\graybold{Preguntas guía:}

\begin{itemize}
\item ¿Debo partir una secuencia o un segmento en trozos, y el costo de cada partición depende del TRAMO completo que se está partiendo?
\item ¿La respuesta de un tramo se arma eligiendo dónde hacer el primer (o último) corte y sumando las respuestas de los dos sub-tramos?
\item ¿Los sub-tramos son siempre intervalos contiguos, de modo que el estado es un par (inicio, fin)?
\item ¿El costo del tramo NO depende de dónde se corte, solo de sus extremos? (es la forma a la que se aplica la optimización de Knuth, si además el costo cumple la desigualdad cuadrangular)
\end{itemize}

\graybold{Utilidad:} Es la plantilla del DP sobre intervalos: el estado es un segmento y la transición prueba todos los puntos interiores donde partirlo. Resuelve multiplicación óptima de cadenas de matrices, corte óptimo de barras o cuerdas, partición en palíndromos, fusión de pilas de piedras y árboles binarios de búsqueda óptimos. Lo que cambia entre esos problemas es el costo de cada partición (en la barra depende solo del tramo; en la cadena de matrices, también del punto de corte); la estructura del recorrido es idéntica.

\graybold{Complejidad:} \bigO{n^3} en la versión directa, o \bigO{n^2} con la optimización de Knuth de los extras.

\graybold{Fundamento teórico:} Se agregan centinelas en los extremos, de modo que los puntos de corte quedan en un arreglo $p[0] = 0 < p_1 < \cdots < p_n < p[n+1] = L$. Sea $dp[i][j]$ el costo mínimo de hacer todos los cortes estrictamente entre $p[i]$ y $p[j]$. La observación central es que, sin importar en qué orden se corte, el PRIMER corte de ese tramo cuesta siempre su longitud completa $p[j] - p[i]$, porque en ese momento el tramo todavía está entero. Elegido ese primer corte en $k$, el problema se separa en dos subproblemas independientes, los tramos $[i,k]$ y $[k,j]$, que nunca vuelven a interactuar. Por lo tanto
\[dp[i][j] = (p[j] - p[i]) + \min_{i<k<j}\big(dp[i][k] + dp[k][j]\big),\]
con $dp[i][i+1] = 0$ porque un tramo sin cortes interiores no cuesta nada. El orden de cálculo es por LONGITUD de intervalo creciente: al resolver un tramo, sus dos sub-tramos son más cortos y ya están listos. La optimización de Knuth se apoya en dos propiedades del costo $w(i,j) = p[j] - p[i]$: la desigualdad cuadrangular, $w(a,c) + w(b,d) \le w(a,d) + w(b,c)$ para $a \le b \le c \le d$ (aquí se cumple con igualdad), y la monotonía, $w(b,c) \le w(a,d)$ cuando $[b,c]$ está dentro de $[a,d]$. Juntas garantizan que el punto de corte óptimo es MONÓTONO: si $\mathrm{opt}[i][j]$ es el mejor $k$, entonces $\mathrm{opt}[i][j-1] \le \mathrm{opt}[i][j] \le \mathrm{opt}[i+1][j]$. Recorriendo solo ese rango, la suma total de iteraciones telescopia y el algoritmo baja a \bigO{n^2}.~\cite{knuth1971}

\graybold{Ejercicio aplicado}

(UVA 10003 — Cutting Sticks) Dada una barra de longitud $l$ y $n$ posiciones donde hay que cortarla, donde cada corte cuesta la longitud del trozo que se está cortando, hallar el costo mínimo total. La entrada termina con $l = 0$.

\graybold{I/O:} Varios casos con centinela $l = 0$ (patrón 3), leídos con lectura total + tokenización porque son números sueltos repartidos en varias líneas. $l < 1000$ y $n < 50$, así que la tabla es de a lo sumo $52\times52$. Verificado contra los dos casos de ejemplo y contrastado con una fuerza bruta que simula TODOS los órdenes de corte posibles (la definición literal del enunciado) en 254 casos con hasta 7 cortes, incluyendo $n = 0$, sin discrepancias. El enunciado no acota la cantidad de casos de prueba, así que se midió el costo por caso con $n = 49$: \aprox{}1.14 ms por caso, es decir \aprox{}1.1s con 1000 casos y \aprox{}5.7s con 5000. La versión con la optimización de Knuth de los extras, verificada igual contra la fuerza bruta y contra esta, da \aprox{}0.35 ms por caso (\aprox{}1.75s con 5000), unas 3.3 veces más rápido: conviene usarla si la entrada real resulta ser grande.

\graybold{Código solución}

\begin{lstlisting}[language=Python]
import sys

def solve():
    data = sys.stdin.buffer.read().split()
    pos = 0
    out = []
    while pos < len(data):
        l = int(data[pos]); pos += 1
        if l == 0:
            break
        n = int(data[pos]); pos += 1
        p = [0] + [int(v) for v in data[pos:pos+n]] + [l]   # centinelas en los extremos
        pos += n
        m = n + 2
        INF = float('inf')
        dp = [[0]*m for _ in range(m)]       # dp[i][j]: cortar lo que hay entre p[i] y p[j]
        # se resuelve por tramos cada vez mas largos: los sub-tramos ya estan listos
        for salto in range(2, m):
            for i in range(m - salto):
                j = i + salto
                largo = p[j] - p[i]          # el PRIMER corte del tramo siempre cuesta esto
                mejor = INF
                for k in range(i + 1, j):    # donde hacer ese primer corte
                    v = dp[i][k] + dp[k][j]
                    if v < mejor: mejor = v
                dp[i][j] = mejor + largo
        out.append("The minimum cutting is %d." % dp[0][m-1])
    sys.stdout.write("\n".join(out) + "\n")

solve()
\end{lstlisting}

\graybold{Extras: optimización de Knuth} — reemplaza el bucle interno completo por un rango acotado, bajando de \bigO{n^3} a \bigO{n^2}. Aplica cuando el costo cumple la desigualdad cuadrangular y es monótono (ver el fundamento), condiciones que aquí se cumplen; si no se cumplen, puede dar respuestas incorrectas, así que conviene verificarla contra la versión cúbica antes de confiar en ella.

\begin{lstlisting}[language=Python]
        dp = [[0]*m for _ in range(m)]
        opt = [[0]*m for _ in range(m)]      # mejor punto de corte de cada tramo
        for i in range(m - 1):
            opt[i][i+1] = i + 1
        for salto in range(2, m):
            for i in range(m - salto):
                j = i + salto
                # el corte optimo es monotono: basta mirar entre estos dos
                lo = opt[i][j-1]; hi = opt[i+1][j]
                mejor = None; arg = lo
                fila = dp[i]
                for k in range(lo, hi + 1):
                    v = fila[k] + dp[k][j]
                    if mejor is None or v < mejor:
                        mejor = v; arg = k
                dp[i][j] = mejor + p[j] - p[i]
                opt[i][j] = arg
\end{lstlisting}
